\documentclass[../../main.tex]{subfiles}

\graphicspath{{../../image/}{../../figure/}}

\begin{document}

\section{习题}\label{section:ch06-习题}

\begin{exercise}\label{exercise:ch06-01}
（★★★）使用极坐标下的 Gauss 积分推导 $D$ 维空间中单位半径超球体的表面积 $S_D$ 和体积 $V_D$ 的计算表达式。为此，请考虑以下结果，该结果是通过从笛卡儿坐标转换为极坐标获得的：
\begin{align}
\prod_{i=1}^{D}\int_{-\infty}^{\infty}e^{-x_i^{2}}\,\mathrm{d}x_i=S_D\int_0^{\infty}e^{-r^{2}}r^{D-1}\,\mathrm{d}r .
\label{align:ch06-Gauss极坐标}
\end{align}
使用 Gamma 函数，定义如下：
\begin{align}
\Gamma(x)=\int_0^{\infty}t^{x-1}e^{-t}\,\mathrm{d}t,
\label{align:ch06-Gamma函数}
\end{align}
以及标准 Gauss 积分 $\int_{-\infty}^{\infty}e^{-x^{2}}\mathrm{d}x=\sqrt{\pi}$，计算\eqref{align:ch06-Gauss极坐标}的两边，从而证明
\begin{align}
S_D=\frac{2\pi^{D/2}}{\Gamma(D/2)} .
\label{align:ch06-超球表面积}
\end{align}
接下来，通过对从 $0$ 到 $1$ 的半径进行积分，证明 $D$ 维空间中单位超球体的体积由下式给出：
\begin{align}
V_D=\frac{S_D}{D} .
\label{align:ch06-超球体积}
\end{align}
最后，使用结果 $\Gamma(1)=1$ 和 $\Gamma(3/2)=\sqrt{\pi}/2$ 证明\eqref{align:ch06-超球表面积}和\eqref{align:ch06-超球体积}可以简化为 $D=2$ 和 $D=3$ 的常用表达式。
\end{exercise}

\begin{exercise}\label{exercise:ch06-02}
（★★★）考虑 $D$ 维空间中半径为 $a$ 的超球体与边长为 $2a$ 的同心超立方体，使超球体在其每条边的中心与超立方体相交。通过使用\cref{exercise:ch06-01}的结果，证明超球体的体积与超立方体的体积之比由下式给出：
\begin{align}
\frac{\text{超球体的体积}}{\text{超立方体的体积}}=\frac{\pi^{D/2}}{D\,2^{D-1}\,\Gamma(D/2)} .
\label{align:ch06-球与立方体体积比}
\end{align}
以如下形式使用Stirling公式：
\begin{align}
\Gamma(x+1)\approx(2\pi)^{1/2}e^{-x}x^{x+1/2},
\label{align:ch06-Stirling}
\end{align}
它对 $x\gg1$ 有效，证明当 $D\to\infty$ 时，\eqref{align:ch06-球与立方体体积比}等于零。同时证明从超立方体的中心到其中一个角的距离，与从超立方体的中心到其中一条边的垂直距离之比为 $\sqrt{D}$，且在 $D\to\infty$ 时趋向于 $\infty$。从这些结果中可以看到，在高维空间中，立方体的大部分体积都集中在大量的角上，这些角本身变成了很长的“尖刺”！
\end{exercise}

\begin{exercise}\label{exercise:ch06-03}
（★★★）探讨 Gauss 分布在高维空间中的行为。考虑 $D$ 维空间中的 Gauss 分布，它由下式给出：
\begin{align}
p(\boldsymbol{x})=\frac{1}{(2\pi\sigma^{2})^{D/2}}\exp\!\left(-\frac{\|\boldsymbol{x}\|^{2}}{2\sigma^{2}}\right).
\label{align:ch06-各向同性Gauss}
\end{align}
我们希望在极坐标中找到以半径为参数的密度函数，其中方向变量已被积分消去。为此，证明半径为 $r$ 且厚度为 $\varepsilon$（$\varepsilon\ll1$）的薄壳上的概率质量由 $p(r)\varepsilon$ 给出，其中
\begin{align}
p(r)=\frac{S_Dr^{D-1}}{(2\pi\sigma^{2})^{D/2}}\exp\!\left(-\frac{r^{2}}{2\sigma^{2}}\right),
\label{align:ch06-径向密度}
\end{align}
$S_D$ 是 $D$ 维空间中单位超球体的表面积。对于较大的 $D$ 值，证明函数 $p(r)$ 有单个位于 $\hat r\approx\sqrt{D}\sigma$ 的驻点。通过考虑 $p(\hat r+\varepsilon)$，其中 $\varepsilon\ll\hat r$，证明对于较大的 $D$ 值，有
\begin{align}
p(\hat r+\varepsilon)\approx p(\hat r)\exp\!\left(-\frac{\varepsilon^{2}}{\sigma^{2}}\right).
\label{align:ch06-径向展开}
\end{align}
这表明 $\hat r$ 是径向概率密度的最大值，且 $p(r)$ 在长度尺度为 $\sigma$ 的最大值 $\hat r$ 处开始指数衰减。对于较大的 $D$ 值来说，$\sigma\ll\hat r$，因此我们看到大部分概率质量集中在半径较大的薄壳内。最后，证明原点处的概率密度 $p(\boldsymbol{x})$ 大于半径 $\hat r$ 处的概率密度，它们之间相差一个因子 $\exp(D/2)$。于是，我们可以看到在高维空间中，Gauss 分布中的大多数概率质量位于与高概率密度区域不同的半径处。
\end{exercise}

\begin{exercise}\label{exercise:ch06-04}
（★★）考虑\eqref{align:ch06-网络函数}给出的网络函数，其中隐藏单元非线性激活函数 $h(\cdot)$ 由 sigmoid 函数给出：
\begin{align}
\sigma(a)=\bigl\{1+\exp(-a)\bigr\}^{-1}.
\label{align:ch06-sigmoid重述}
\end{align}
证明存在一个等价网络，它计算的是完全相同的网络函数，但具有由 $\tanh(a)$ 给出的隐藏单元激活函数，其中 tanh 函数由\eqref{align:ch06-tanh}定义。提示：首先找到 $\sigma(a)$ 和 $\tanh(a)$ 之间的关系，然后证明这两个网络的参数可以通过线性变换相互转换。
\end{exercise}

\begin{exercise}\label{exercise:ch06-05}
（★★）swish 激活函数（Ramachandran，Zoph，and Le，2017）定义如下：
\begin{align}
h(x)=x\,\sigma(\beta x),
\label{align:ch06-swish}
\end{align}
其中 $\sigma(x)$ 是\eqref{align:ch06-sigmoid}定义的 sigmoid 激活函数。当在神经网络中使用时，$\beta$ 可以视为可学习参数。手绘或用软件画出 swish 激活函数及其在 $\beta=0.1$、$\beta=1.0$ 和 $\beta=10$ 时的一阶导数。证明当 $\beta\to\infty$ 时，swish 函数变为 ReLU 函数。
\end{exercise}

\begin{exercise}\label{exercise:ch06-06}
（★）我们在 sigmoid 激活函数的讨论中看到，其导数可以用函数值本身来表示。推导\eqref{align:ch06-tanh}定义的 tanh 激活函数的相应结果。
\end{exercise}

\begin{exercise}\label{exercise:ch06-07}
（★★）证明\eqref{align:ch06-softplus}给出的 softplus 激活函数 $\zeta(a)$ 满足以下性质：
\begin{align}
\zeta(a)-\zeta(-a)&=a,
\label{align:ch06-softplus奇}
\\
\ln\sigma(a)&=-\zeta(-a),
\label{align:ch06-softplus对数}
\\
\frac{\mathrm{d}\zeta(a)}{\mathrm{d}a}&=\sigma(a),
\label{align:ch06-softplus导数}
\\
\zeta^{-1}(a)&=\ln\bigl(\exp(a)-1\bigr),
\label{align:ch06-softplus逆}
\end{align}
其中 $\sigma(a)$ 是\eqref{align:ch06-sigmoid}给出的 sigmoid 激活函数。
\end{exercise}

\begin{exercise}\label{exercise:ch06-08}
（★）证明对方差 $\sigma^{2}$ 最小化误差函数［\eqref{align:ch06-回归负对数似然}］可以得到\eqref{align:ch06-噪声方差}。
\end{exercise}

\begin{exercise}\label{exercise:ch06-09}
（★）证明在多输出神经网络的条件分布［\eqref{align:ch06-多输出条件Gauss}］下最大化似然函数等价于最小化平方和误差函数［\eqref{align:ch06-多输出平方和}］。此外，证明最小化平方和误差函数的噪声方差由\eqref{align:ch06-多输出噪声方差}给出。
\end{exercise}

\begin{exercise}\label{exercise:ch06-10}
（★★）考虑一个涉及多个目标变量的回归问题，其中假设以输入向量 $\boldsymbol{x}$ 为条件的目标分布是具有以下形式的 Gauss 分布：
\begin{align}
p(\boldsymbol{t}\mid\boldsymbol{x},\boldsymbol{w})=\mathcal{N}\bigl(\boldsymbol{t}\mid\boldsymbol{y}(\boldsymbol{x},\boldsymbol{w}),\boldsymbol{\Sigma}\bigr),
\label{align:ch06-一般协方差Gauss}
\end{align}
其中 $\boldsymbol{y}(\boldsymbol{x},\boldsymbol{w})$ 是具有输入向量 $\boldsymbol{x}$ 和权重向量 $\boldsymbol{w}$ 的神经网络的输出，$\boldsymbol{\Sigma}$ 是假定的目标上 Gauss 噪声的协方差。给定一组 $\boldsymbol{x}$ 和 $\boldsymbol{t}$ 的独立观测值，假设 $\boldsymbol{\Sigma}$ 是固定的且已知，写出为找到 $\boldsymbol{w}$ 的最大似然解而必须最小化的误差函数。假设 $\boldsymbol{\Sigma}$ 也要从数据中确定，写出 $\boldsymbol{\Sigma}$ 最大似然解的表达式。注意，现在 $\boldsymbol{w}$ 和 $\boldsymbol{\Sigma}$ 的优化是耦合的，这与\cref{subsection:ch06-回归}讨论的独立目标变量情况相反。
\end{exercise}

\begin{exercise}\label{exercise:ch06-11}
（★★）考虑一个二分类问题，其中目标值为 $t\in\{0,1\}$，用网络输出 $y(\boldsymbol{x},\boldsymbol{w})$ 表示 $p(t=1\mid\boldsymbol{x})$，并假设训练数据点上的类标签设置不正确的概率为 $\varepsilon$。假设数据独立同分布，写出对应于负对数似然的误差函数。验证当 $\varepsilon=0$ 时将得到误差函数［\eqref{align:ch06-二分类交叉熵}］。注意，与交叉熵误差函数相比，此误差函数使得模型对错误标记的数据具有鲁棒性。
\end{exercise}

\begin{exercise}\label{exercise:ch06-12}
（★★）二分类问题的误差函数［\eqref{align:ch06-二分类交叉熵}］是针对具有 sigmoid 输出单元激活函数的网络推导出来的，故有 $0\leqslant y(\boldsymbol{x},\boldsymbol{w})\leqslant1$，且目标值为 $t\in\{0,1\}$。考虑一个具有输出 $-1\leqslant y(\boldsymbol{x},\boldsymbol{w})\leqslant1$ 的网络，对类 $\mathcal{C}_1$ 目标值为 $t=1$，而对类 $\mathcal{C}_2$ 目标值为 $t=-1$，推导出相应的误差函数。进一步地，输出单元激活函数应该如何选择？
\end{exercise}

\begin{exercise}\label{exercise:ch06-13}
（★）证明最大化多分类神经网络模型的似然，其中网络输出可解释为 $y_k(\boldsymbol{x},\boldsymbol{w})=p(t_k=1\mid\boldsymbol{x})$，等价于最小化交叉熵误差函数［\eqref{align:ch06-多分类交叉熵}］。
\end{exercise}

\begin{exercise}\label{exercise:ch06-14}
（★）证明对于具有 sigmoid 激活函数 $y_k=\sigma(a_k)$ 的输出单元，误差函数［\eqref{align:ch06-二分类交叉熵}］对其预激活 $a_k$ 的导数满足\eqref{align:ch06-输出单元导数}，其中 $\sigma(a)$ 由\eqref{align:ch06-sigmoid}给出。
\end{exercise}

\begin{exercise}\label{exercise:ch06-15}
（★）证明对于具有 softmax 激活函数［\eqref{align:ch06-softmax}］的输出单元，误差函数［\eqref{align:ch06-多分类交叉熵}］对其预激活 $a_k$ 的导数满足\eqref{align:ch06-输出单元导数}。
\end{exercise}

\begin{exercise}\label{exercise:ch06-16}
（★★）用一对关于关节角 $\theta_1$ 和 $\theta_2$ 以及连接长度 $L_1$ 和 $L_2$ 的方程表示\cref{figure:ch06-机械臂}所示机械臂的笛卡儿坐标 $(x_1,x_2)$。假设坐标系的原点由下臂的连接点给出。这些方程定义了机械臂的正运动学。
\end{exercise}

\begin{exercise}\label{exercise:ch06-17}
（★★）证明\eqref{align:ch06-分量后验}定义的变量 $\gamma_{nk}$ 可以看作混合分布［\eqref{align:ch06-混合密度模型}］分量的后验概率 $p(k\mid\boldsymbol{t})$，其中混合系数 $\pi_k(\boldsymbol{x})$ 可以视为与 $\boldsymbol{x}$ 相关的先验概率 $p(k)$。
\end{exercise}

\begin{exercise}\label{exercise:ch06-18}
（★★）对于控制混合密度网络中混合系数的网络输出预激活，推导误差函数对其的导数［\eqref{align:ch06-混合系数导数}］。
\end{exercise}

\begin{exercise}\label{exercise:ch06-19}
（★★）对于控制混合密度网络中分量均值的网络输出预激活，推导误差函数对其的导数［\eqref{align:ch06-均值导数}］。
\end{exercise}

\begin{exercise}\label{exercise:ch06-20}
（★★）对于控制混合密度网络中分量方差的网络输出预激活，推导误差函数对其的导数［\eqref{align:ch06-方差导数}］。
\end{exercise}

\begin{exercise}\label{exercise:ch06-21}
（★★★）针对混合密度网络模型的条件均值和方差，验证\eqref{align:ch06-条件均值}和\eqref{align:ch06-条件方差}给出的结果。
\end{exercise}

\end{document}
